\uuid{548P}
\titre{Intervalle centré de probabilité 0,95 pour une loi normale}
\auteur{Maxime Nguyen}
\organisation{AMSCC}
\datecreate{2026-11-27}
\niveau{L2}
\module{Probabilité}
\chapitre{Probabilité continue}
\sousChapitre{Loi normale}
\difficulte{2}
\theme{loi normale, intervalle centré, quantile, intervalle de confiance}

\contenu{
\begin{question}
\qtype{onlyone}
On note $\mathcal{N}(\mu,\sigma)$ la loi normale d'espérance $\mu$ et d'écart-type $\sigma$, et $\Phi$ la fonction de répartition de $\mathcal{N}(0,1)$. On donne $\Phi(1{,}645) \approx 0{,}95$ et $\Phi(1{,}96) \approx 0{,}975$. Soit $X \sim \mathcal{N}(\mu, \sigma)$. Le réel $h > 0$ tel que $\mathbb{P}(\mu - h \leq X \leq \mu + h) = 0{,}95$ vaut environ :
\begin{answers}
\bad{$1{,}96$}
\bad{$1{,}645\,\sigma$}
\bad{$1{,}96\,\sigma^2$}
\good{$1{,}96\,\sigma$}
\end{answers}
\begin{explanations}
En posant $t = \frac{h}{\sigma}$, on veut
$\mathbb{P}(-t \leq Z \leq t) = 2\Phi(t) - 1 = 0{,}95$, soit $\Phi(t) = 0{,}975$,
donc $t \approx 1{,}96$ et $h \approx 1{,}96\,\sigma$.

A oublie de revenir à l'échelle de $X$.

B : $\Phi(1{,}645) = 0{,}95$ laisse $5\,\%$ dans une queue ; un intervalle centré à $95\,\%$ laisse $2{,}5\,\%$ de chaque côté.

C multiplie par la variance.

C'est le calcul qui sert, au chapitre sur l'estimation, à construire les
intervalles de confiance à $95\,\%$.

Contre-sens visé : chercher $\Phi(t) = 0{,}95$ au lieu de
$\Phi(t) = 0{,}975$, c'est-à-dire oublier qu'un intervalle centré partage le
risque entre deux queues.
\end{explanations}
\end{question}
}
